% Threshold corollaries for Bernoulli(p) baseline (kept separate to avoid inflating the main baseline file).

\begin{corollary}[单点输入--输出位的唯一独立阈值]\label{cor:fold-gauge-anomaly-bernoulli-p-bitpair-independence-threshold}
在命题 \ref{prop:fold-gauge-anomaly-bernoulli-p-bitpair-law} 的平稳 Bernoulli\((p)\) 基线设定下，
令 \(p_\phi:=\frac{\sqrt5-1}{2}\in(0,1)\)。
则随机变量 \(X_0,Y_0\in\{0,1\}\) 严格独立当且仅当 \(p=p_\phi\)。
\leanverified{paper\_fold\_gauge\_anomaly\_bitpair\_independence\_threshold}
\end{corollary}
\begin{proof}
由命题 \ref{prop:fold-gauge-anomaly-bernoulli-p-bitpair-law}，
\[
\Cov(X_0,Y_0)
=P(X_0=1,Y_0=1)-P(X_0=1)P(Y_0=1)
=-\frac{p(1-p)^2(p^2+p-1)}{1+p^3}.
\]
对 \(p\in(0,1)\) 有 \(p(1-p)^2/(1+p^3)>0\)，故 \(\Cov(X_0,Y_0)=0\) 当且仅当 \(p^2+p-1=0\)；
其在 \((0,1)\) 的唯一根为 \(p_\phi\)。
由于 \(\{0,1\}\)-值随机变量的联合分布由边缘分布与 \(P(X_0=1,Y_0=1)\) 唯一确定，
\(\Cov(X_0,Y_0)=0\) 等价于 \(P(X_0=1,Y_0=1)=P(X_0=1)P(Y_0=1)\)，从而等价于独立性。
证毕。
\end{proof}

\begin{corollary}[失配密度极大点的 Cardano 闭式]\label{cor:fold-gauge-anomaly-bernoulli-p-density-maximizer-cardano}
令 \(g_\ast(p)=\frac{p^2(3-2p)}{1+p^3}\)。
记 \(p_\star\in(0,1)\) 为三次方程
\[
p^3+2p-2=0
\]
在 \((0,1)\) 中的唯一实根（推论 \ref{cor:fold-gauge-anomaly-bernoulli-p-density-unimodal}）。
则该根可由 Cardano 公式写为
\[
p_\star=\sqrt[3]{\,1+\sqrt{\tfrac{35}{27}}\,}+\sqrt[3]{\,1-\sqrt{\tfrac{35}{27}}\,},
\]
并且极大值满足恒等式 \(g_\ast(p_\star)=p_\star^2\)。
\leanverified{paper\_fold\_gauge\_anomaly\_density\_maximizer\_cardano}
\end{corollary}
\begin{proof}
由推论 \ref{cor:fold-gauge-anomaly-bernoulli-p-density-unimodal}，\(p_\star\) 为所示三次方程在 \((0,1)\) 的唯一根，
且 \(g_\ast(p_\star)=p_\star^2\) 已在该推论中给出。
将方程写为降次形式 \(p^3+2p=2\)，应用 Cardano 公式得到所示闭式表示。证毕。
\end{proof}

\begin{theorem}[Gallavotti--Cohen 缺陷的线性响应阈值]\label{thm:fold-gauge-anomaly-gc-defect-linear-response-threshold}
在定理 \ref{thm:fold-gauge-anomaly-bernoulli-p-laws} 的 Bernoulli\((p)\) 基线设定下，定义压力
\[
P_{G,p}(\theta)=\log\lambda(e^\theta,p),
\]
以及 Gallavotti--Cohen 缺陷函数
\[
\Delta_p(\theta):=P_{G,p}(\theta)-\theta-P_{G,p}(-\theta).
\]
则 \(\Delta_p\) 在 \(\theta=0\) 的泰勒展开满足
\[
\Delta_p(\theta)=\bigl(2g_\ast(p)-1\bigr)\theta+O(\theta^3),
\qquad \theta\to 0,
\]
其中失配密度 \(g_\ast(p)=\frac{p^2(3-2p)}{1+p^3}\)。
因此存在唯一阈值
\[
p_0=\frac{1+\sqrt{21}}{10}\in(0,1)
\]
使得当 \(p<p_0\) 时，\(\Delta_p(\theta)<0\) 对一切充分小的 \(\theta>0\) 成立；当 \(p>p_0\) 时，\(\Delta_p(\theta)>0\) 对一切充分小的 \(\theta>0\) 成立。
\leanverified{paper\_fold\_gauge\_anomaly\_gc\_defect\_linear\_response\_threshold}
\end{theorem}

\begin{proof}
由定理 \ref{thm:fold-gauge-anomaly-bernoulli-p-laws}(4)，\(P_{G,p}\) 在 \(\theta=0\) 的邻域实解析，并且
\(
P_{G,p}'(0)=g_\ast(p)
\)。
由定义，
\[
\Delta_p'(0)=P_{G,p}'(0)-1-(-P_{G,p}'(0))=2P_{G,p}'(0)-1=2g_\ast(p)-1.
\]
由于 \(\Delta_p\) 为奇函数（见命题 \ref{prop:fold-gauge-anomaly-bernoulli-p-gc-defect-farfield-critical} 的证明开头），其泰勒展开只含奇次项，从而
\(
\Delta_p(\theta)=\Delta_p'(0)\theta+O(\theta^3)
\)。
代入 \(g_\ast(p)=\frac{p^2(3-2p)}{1+p^3}\) 得
\[
\Delta_p'(0)=\frac{6p^2-5p^3-1}{1+p^3}
=\frac{(1-p)(5p^2-p-1)}{1+p^3}.
\]
在区间 \(p\in(0,1)\) 上，分母与因子 \(1-p\) 严格为正，故 \(\Delta_p'(0)\) 的符号由二次多项式 \(5p^2-p-1\) 的符号决定。
该二次方程在 \((0,1)\) 内存在且仅存在一根
\(
p_0=\frac{1+\sqrt{21}}{10}
\)。
结合 \(\Delta_p(\theta)=\Delta_p'(0)\theta+O(\theta^3)\) 即得小正 \(\theta\) 下的符号结论。证毕。
\end{proof}

\begin{proposition}[高场一阶指数渐近与黄金点闭式系数]\label{prop:fold-gauge-anomaly-gc-defect-farfield-expansion}
\leanverified{paper\_fold\_gauge\_anomaly\_gc\_defect\_farfield\_expansion}
在命题 \ref{prop:fold-gauge-anomaly-bernoulli-p-gc-defect-farfield-critical} 的记号下，令
\[
\alpha(p):=\frac{1-p+\sqrt{(1-p)(1+3p)}}{2},
\qquad
c(p):=\bigl(p^2(1-p)\bigr)^{1/3}.
\]
则存在显式函数 \(A(p)\) 使当 \(\theta\to+\infty\) 时
\[
\Delta_p(\theta)=\log\frac{c(p)}{\alpha(p)}+A(p)e^{-\theta}+O(e^{-2\theta}),
\]
并且可以取
\[
A(p)=\frac{d(p)}{c(p)},
\qquad
d(p)=\frac13\left(\left(\frac{p}{1-p}\right)^{1/3}+(1-p)\right).
\]
特别地，在黄金偏置点 \(p=\frac{1+\sqrt5}{4}\) 处有 \(\log\frac{c(p)}{\alpha(p)}=0\) 且
\[
\Delta_{p}(\theta)=\frac{5+\sqrt5}{6}e^{-\theta}+O(e^{-2\theta}),
\qquad \theta\to+\infty.
\]
\end{proposition}

\begin{proof}
令 \(u=e^\theta\) 并记 \(\lambda(u,p)\) 为定理 \ref{thm:fold-gauge-anomaly-bernoulli-p-laws} 中的 Perron 根（满足四次方程 \eqref{eq:fold-gauge-anomaly-bernoulli-p-pressure-quartic}）。
由定义
\[
\Delta_p(\theta)=\log\frac{\lambda(u,p)}{u\,\lambda(u^{-1},p)}.
\]
\par
\textbf{(i) 大场端 \(u\to\infty\)。}
设 \(\lambda(u,p)=u\,s(u,p)\)。
将其代入 \eqref{eq:fold-gauge-anomaly-bernoulli-p-pressure-quartic} 并以 \(u^4\) 归一化，可得
\[
s(u,p)^4-p^2(1-p)\,s(u,p)+O(u^{-1})=0,
\qquad u\to\infty.
\]
由 Perron 分支的正性与唯一性，\(s(u,p)\to c(p)\) 且满足 \(c(p)^3=p^2(1-p)\)。
进一步令
\(
\lambda(u,p)=c(p)u+d+O(u^{-1})
\)。
将该展开代回 \eqref{eq:fold-gauge-anomaly-bernoulli-p-pressure-quartic}，最高阶 \(u^4\) 项由 \(c(p)^3=p^2(1-p)\) 消去；比较 \(u^3\) 系数得到一条线性方程，从而唯一确定
\[
d(p)=\frac13\left(\left(\frac{p}{1-p}\right)^{1/3}+(1-p)\right).
\]
因此
\[
\log\frac{\lambda(u,p)}{u}
=\log\left(c(p)+\frac{d(p)}{u}+O(u^{-2})\right)
=\log c(p)+\frac{d(p)}{c(p)}u^{-1}+O(u^{-2}).
\]
\par
\textbf{(ii) 小场端 \(u\to 0^+\)。}
由命题 \ref{prop:fold-gauge-anomaly-bernoulli-p-endpoint-ldp}，有 \(\lambda(0,p)=\alpha(p)\)。
对四次方程 \eqref{eq:fold-gauge-anomaly-bernoulli-p-pressure-quartic} 记其左端为 \(F(\lambda,u,p)\)。
在 \((\lambda,u)=(\alpha(p),0)\) 处有 \(F(\alpha,0,p)=0\)。
利用 \(\alpha(p)\) 满足二次关系 \(\alpha^2-(1-p)\alpha-p(1-p)=0\)，可直接化简得到
\[
\partial_u F(\alpha(p),0,p)=0,
\]
从而 Perron 分支在 \(u=0\) 处的一阶项消失，并存在常数 \(C(p)<\infty\) 使
\[
\lambda(u,p)=\alpha(p)+O(u^2),
\qquad u\to 0^+.
\]
于是
\[
\log\lambda(u^{-1},p)=\log\alpha(p)+O(u^{-2}),
\qquad u\to\infty.
\]
\par
将 (i)(ii) 合并并取 \(u=e^\theta\) 即得
\[
\Delta_p(\theta)=\log\frac{c(p)}{\alpha(p)}+\frac{d(p)}{c(p)}e^{-\theta}+O(e^{-2\theta}).
\]
黄金偏置点由命题 \ref{prop:fold-gauge-anomaly-bernoulli-p-gc-defect-farfield-critical} 给出为 \(p=\frac{1+\sqrt5}{4}\)，此时 \(\log\frac{c(p)}{\alpha(p)}=0\) 且 \(p/(1-p)=\varphi^3\)，从而
\(
d(p)=\frac13(\varphi+(1-p))
\)、
\(
c(p)=\tfrac12
\)，化简得到 \(A(p)=d(p)/c(p)=\frac{5+\sqrt5}{6}\)。
证毕。
\end{proof}

\begin{corollary}[必然零点区间与区间夹逼]\label{cor:fold-gauge-anomaly-gc-defect-must-cross}
\leanverified{paper\_fold\_gauge\_anomaly\_gc\_defect\_must\_cross}
令
\(
p_0=\frac{1+\sqrt{21}}{10}
\)（定理 \ref{thm:fold-gauge-anomaly-gc-defect-linear-response-threshold}）
与
\(
p_\infty=\frac{1+\sqrt5}{4}
\)（命题 \ref{prop:fold-gauge-anomaly-bernoulli-p-gc-defect-farfield-critical}）。
则对任意 \(p\in(p_0,p_\infty)\)，存在 \(\theta_p>0\) 使得 \(\Delta_p(\theta_p)=0\)。
\end{corollary}

\begin{proof}
由定理 \ref{thm:fold-gauge-anomaly-gc-defect-linear-response-threshold}，当 \(p>p_0\) 时，存在 \(\varepsilon>0\) 使 \(\Delta_p(\theta)>0\) 对所有 \(\theta\in(0,\varepsilon)\) 成立。
另一方面，由命题 \ref{prop:fold-gauge-anomaly-bernoulli-p-gc-defect-farfield-critical}，当 \(p<p_\infty\) 时有
\(
\lim_{\theta\to+\infty}\Delta_p(\theta)<0
\)，故存在 \(T>0\) 使得对所有 \(\theta>T\) 有 \(\Delta_p(\theta)<0\)。
由于 \(\Delta_p\) 在 \((0,\infty)\) 上连续，介值定理给出某个 \(\theta_p\in(\varepsilon,T)\) 使 \(\Delta_p(\theta_p)=0\)。证毕。
\end{proof}

\begin{corollary}[黄金偏置处的零温中性与一阶指数衰减]\label{cor:fold-gauge-anomaly-gc-defect-golden-neutral}
\leanverified{paper\_fold\_gauge\_anomaly\_gc\_defect\_golden\_neutral}
取 \(p=\frac{1+\sqrt5}{4}\)。
则 \(\Delta_p(\theta)\to 0\) 且满足一阶指数衰减
\[
\Delta_p(\theta)=\frac{5+\sqrt5}{6}e^{-\theta}+O(e^{-2\theta}),
\qquad \theta\to+\infty.
\]
\end{corollary}

\begin{proof}
由命题 \ref{prop:fold-gauge-anomaly-gc-defect-farfield-expansion} 的特例直接得到。证毕。
\end{proof}

\endinput
